\documentclass[10pt,a4paper]{amsart}
\usepackage[margin=19mm]{geometry}
\usepackage{amsmath,amssymb,mathtools}
\usepackage[scr=rsfs,cal=euler]{mathalfa}
\usepackage{xfrac}
\usepackage[unicode,hidelinks,bookmarks=false]{hyperref}
\newcommand{\ie}{i.e.\ }
\newcommand{\cf}{cf.\ }
\newcommand{\resp}{respectively\ }
\newcommand{\Subsection}{\subsection}
\providecommand{\nobreakdash}{\nobreak\hbox{-}}
\setlength{\emergencystretch}{2em}
\multlinegap0pt
\usepackage{bm}
\usepackage{bbm}
\usepackage{dsfont}
\usepackage{xspace}
\usepackage{cancel}
\usepackage{mathtools}
\usepackage{amsthm}
\makeatletter
\@ifundefined{theorem}{\newtheorem{theorem}{Theorem}}{}
\@ifundefined{lemma}{\newtheorem{lemma}[theorem]{Lemma}}{}
\makeatother
\makeatletter
\newcommand{\EK}{\mathsf{E}}
\newcommand{\K}{\mathbb{K}}
\newcommand{\V}{\mathbb{V}}
\newcommand{\dd}{\mathsf{d}}
\renewcommand\div{\operatorname{div}}
\newcommand{\dive}{{\ensuremath\mathop{\mathrm{div}\,}}} % space needed!
\DeclareMathOperator{\divg}{div}
\expandafter\ifx\csname myproof\endcsname\relax
\newenvironment{myproof}[1][\proofname]{%
  \proof Proof of #1%
}{\endproof}
\fi
\newcommand{\og}{\omega}
\newcommand\skw{\operatorname{skw}}
\makeatother
\title[The Johnson-Krizek-Mercier elasticity element in any dimensi]{The Johnson-Krizek-Mercier elasticity element in any dimensions}
\author{Jay Gopalakrishnan, Johnny Guzman, Jeonghun J. Lee}
\date{Source: arXiv:2403.13189v4 (CC BY 4.0)}
\begin{document}
\maketitle

\noindent\emph{Source and licence.} The Johnson-Krizek-Mercier elasticity element in any dimensions. Version arXiv:2403.13189v4 (CC BY 4.0), distributed under the Creative Commons Attribution 4.0 International licence, as recorded in the arXiv licence metadata for this exact version and shown on the abstract page https://arxiv.org/abs/2403.13189v4. Full rights notice: licenses/arxiv-2403.13189-CC-BY-4.0.txt.\par\smallskip

\noindent\textit{Editorial scope.} Counted unit: the Lemma whose source label is \texttt{applemma2-1} in the article (source0.tex lines 2398--2413, source label \texttt{applemma2-1}), delivered together with the article's own complete original proof of it (source0.tex lines 2414--2455, one \texttt{proof} environment of 42 lines). No mathematics is written, repaired or re-derived: statement and proof are the article's own text line for line. Required context retained uncounted: eq:div-KV (lines 2268--2271); eq:dd-2 (lines 2288--2293); eq:Theta-def-2 (lines 2392--2396). Every definition, display and label of those lines is retained unchanged. Omitted from this package and not redistributed: 1--2267, 2272--2287, 2294--2391, 2397--2397, 2456--3259. Nothing inside a retained excerpt is omitted or altered: every retained body is delivered exactly as the article's lines are written, so no omission receipt of any kind accompanies this package. Citations: the retained excerpts carry no citation command at all, so no bibliography entry of the article is transcribed and no reference is redistributed. Reference adaptation: none was needed and none was made; every reference inside the delivered unit resolves inside it, so no unresolved reference or citation is produced. Printed slips kept literal: no printed word, symbol, hypothesis or number of the counted statement or of its proof was altered. Layout and class: the article's own class is \texttt{amsart} with packages including geometry, graphicx, amssymb, amsmath, amsbsy, eucal, amsfonts, mathrsfs, amscd, bm, tcolorbox, mathtools, hyperref, xy; the wrapper is the lane's pinned \texttt{amsart} editorial wrapper at 10pt on A4, which restates the theorem-like environments the retained text needs and adds the article's own macro definitions that the retained text needs (\texttt{\textbackslash EK}, \texttt{\textbackslash K}, \texttt{\textbackslash V}, \texttt{\textbackslash dd}, \texttt{\textbackslash div}, \texttt{\textbackslash dive}, \texttt{\textbackslash divg}, \texttt{\textbackslash og}, \texttt{\textbackslash skw}), with extra wrapper packages (bm, bbm, dsfont, xspace, cancel, mathtools). The delivered numbering is the unit's own. Assets: no retained span contains a figure environment, a graphics-inclusion command, a PDF-page inclusion, a table or a TikZ picture; no image file is copied into this package, the package declares no asset dependency and no image file is redistributed with it. Licence: version arXiv:2403.13189v4 of this article is distributed under the Creative Commons Attribution 4.0 International licence, as recorded in the arXiv licence metadata for this exact version and shown on the abstract page https://arxiv.org/abs/2403.13189v4. A full rights notice is delivered at licenses/arxiv-2403.13189-CC-BY-4.0.txt. Characters: every retained excerpt is pure ASCII, so the delivered engine, XeLaTeX with its default Latin Modern text font, reports no missing character.
\begin{equation}
  \label{eq:div-KV}
  \divg (b_{ijk} \EK_{ij} \otimes e_k) = (\partial_k b_{ijk}) \EK_{ij}.
\end{equation}

\begin{equation}
  \label{eq:dd-2}
  \dd ( a_{ijk} e_i \otimes \EK_{jk}) = \partial_k( a_{ijk} - a_{ikj}) \,
  e_i \otimes e_j,
  \qquad  a_{ijk} e_i \otimes \EK_{jk}\in \Lambda(\V \otimes \K).  
\end{equation}

\begin{equation}
  \label{eq:Theta-def-2}
\Theta (a_{ijk} e_i \otimes \EK_{jk}) = 2 a_{ijk} \EK_{ij} \otimes
e_k.  
\end{equation}

\begin{lemma}
  The operator
  $ \Theta: \Lambda(\V \otimes \K) \rightarrow \Lambda(\K \otimes \V)$
  is invertible and  its inverse is given by
  \[
    \Theta^{-1}\, b
    = \frac 1 2 (b_{ijk} - b_{ikj} - b_{jki}) \,e_i \otimes e_j \otimes e_k
  \]
  for any $b = b_{ijk} \, e_i \otimes e_j \otimes e_k\in \Lambda(\K \otimes \V)$.
  Moreover, for all $\omega \in \Lambda(\V \otimes \K)$, 
  \begin{align}
    \dive \Theta \,\omega 
    & = 2 \skw \dd \,\omega. 
      \label{applemma2-1}       
  \end{align}
\end{lemma}

\begin{proof}
  By direct calculation, it is easy to verify that the given
  expression for $\Theta^{-1}$ satisfies $ \Theta (\Theta^{-1} b) = b$
  for all $b \in \Lambda(\K \otimes \V)$ and $ \Theta^{-1} (\Theta a) = a$ for all
  $a \in \Lambda(\V \otimes \K)$. To detail the latter, letting
  $\theta = \Theta a = (a_{ijk} - a_{jik}) e_i \otimes e_j \otimes
  e_k,$ observe that
  \begin{align*}
    2 \Theta^{-1} \Theta a
    &= 
      (\theta_{ijk} - \theta_{ikj} - \theta_{jki}) e_i \otimes e_j \otimes   e_k
    \\
    & =
      ( (a_{ijk} - a_{jik}) - (a_{ikj} - a_{kij}) - (a_{jki} - a_{kji})
      ) e_i \otimes e_j \otimes  e_k
  \end{align*}
  When $a \in \Lambda(\V \otimes \K)$, its components have skew
  symmetry in the last two indices (i.e., $a_{ijk} = -a_{ikj}$), so
  the last expression above simplifies to
  $2a_{ijk} e_i \otimes e_j \otimes e_k = 2a$.

  Next, to prove~\eqref{applemma2-1}, let
  $\og = \omega_{ijk} e_i \otimes \EK_{jk}\in \Lambda( \V \otimes
  \K)$. Then, by~\eqref{eq:div-KV} and \eqref{eq:Theta-def-2}, 
  \begin{align*}
    \div (\Theta \omega)
    &  % = \div( (\omega_{ijk} - \omega_{jik}) e_i \otimes e_j \otimes e_k)
      = \div( 2\omega_{ijk} \EK_{ij} \otimes e_k)
     = 2(\partial_k \omega_{ijk})\, \EK_{ij}.
  \end{align*}
  Also, by~\eqref{eq:dd-2}, 
  \begin{align*}
    2\skw  \dd \omega
    & 
      % \skw  \dd (  \omega_{ijk} e_i \otimes \EK_{jk})
      =  \partial_k (\omega_{ijk} - \omega_{ikj}) \,2 \skw(e_i \otimes e_j)
    \\
    &  = \partial_k (\omega_{ijk} - \omega_{jik}) \; \EK_{ij} 
      = 2(\partial_k \omega_{ijk})\, \EK_{ij}.
  \end{align*}
  Thus the left and right hand sides of~\eqref{applemma2-1} are equal.
\end{proof}

\end{document}
